% !TeX root = ../../Book3.tex
\section{Ahlfors–Schwarz 引理}
\label{b3:sec:1402}

Ahlfors--Schwarz 比较使用一个在边界发散的圆盘密度，将正的最大值强迫到内部，再由 Laplace 不等式排除它。这个证明只需二阶微分与最大值原理。

\subsection{曲率比较}
\label{b3:subsec:140201}

\term[ahlfors-schwarz-yinli]{Ahlfors--Schwarz 引理}[Ahlfors--Schwarz lemma]用统一负曲率界比较密度。

\begin{theorem}[Ahlfors--Schwarz]\label{b3:thm:14-ahlfors-schwarz}
若 \(\rho>0\) 为 \(D(0,R)\) 上 \(C^2\) 密度，\(K_\rho\le-\kappa<0\)，则
\[
 \rho(z)\le\frac{2R}{\sqrt\kappa(R^2-|z|^2)}.
\]
\end{theorem}
\begin{proof}
固定 \(r<R\)，令 \(\eta_r=2r/(\sqrt\kappa(r^2-|z|^2))\)。它满足 \(\Delta\log\eta_r=\kappa\eta_r^2\)。\(\rho\) 在闭半径 \(r\) 盘有界，\(\eta_r\) 在边界发散，所以 \(\rho/\eta_r\rightarrow0\)。若比值曾大于1，\(v=\log(\rho/\eta_r)\) 有内部正最大值。在该点
\[
 \Delta v\ge\kappa(\rho^2-\eta_r^2)>0,
\]
与最大值处 \(\Delta v\le0\) 矛盾。故 \(\rho\le\eta_r\)，令 \(r\uparrow R\) 即得结论。
\end{proof}
固定较小圆盘再令半径增大，是为了只使用内部有界性，不假定原密度有边界延拓。等号密度本身在最外边界无界，并且准确达到比较常数。

\ProseParagraphBreak
比较证明先用半径 \(r<R\) 的模型，是为创造一个可达到的内部最大值。模型密度在 \(|z|=r\) 发散，而给定密度在闭 \(r\) 盘附近有界，所以比值趋于0，正的最大比值只能在内部出现。若直接在半径 \(R\) 的开盘上写“取最大点”，没有保证比值达到上确界，证明便缺少关键一步。最后让 \(r\uparrow R\)，才得到所需的边界极限比较。

\ProseParagraphBreak
曲率界也需是统一界 \(K\le-\kappa<0\)。若只知道每一点 \(K<0\)，不同半径上的最优 \(\kappa\) 可能趋于0，比较式右边随之增大，不能再让盘半径趋于无穷推出密度为0。选读小节的完整负曲率平面例子正说明这一步不能省略。Ahlfors--Schwarz 是定量的曲率比较，而不是仅由曲率符号决定全部全局几何。

\subsection{Schwarz–Pick 推广}
\label{b3:subsec:140202}

\begin{corollary}\label{b3:cor:14-schwarz-negative-target}
设 \(f:\mathbb D\rightarrow Y\) 全纯，目标有正 \(C^2\) 共形密度 \(\rho\) 且曲率不超过 \(-\kappa\)。则
\[
 \rho(f(z))|f'(z)|\le\frac{2}{\sqrt\kappa(1-|z|^2)}.
\]
\end{corollary}
\begin{proof}
若 \(f\) 常值，左端为0。否则其导数零点离散，拉回密度 \(h\) 连续，零点处为0；非零处满足同一曲率不等式。比较 \(h/\eta_r\) 的正最大值若大于1，既不在边界也不在零点，故完全适用前一定理的 Laplace 论证，得到相同上界。
\end{proof}
取目标圆盘 Poincaré 密度、\(\kappa=1\)，恢复 Schwarz--Pick 导数式。曲率方法并不要求目标的显式自同构；只要有合适密度，便可控制所有源映射。

\subsection{完备度量形式}
\label{b3:subsec:140203}

\begin{theorem}\label{b3:thm:14-hyperbolic-density-maximal}
若 \(X\) 为双曲曲面，\(\rho\) 是其任一正 \(C^2\) 共形密度且曲率不超过 \(-\kappa\)，则 \(\rho\le\lambda_X/\sqrt\kappa\)。因此对 \(f:X\rightarrow Y\) 及前一推论的目标密度，有 \(\rho(f)|f'|\le\lambda_X/\sqrt\kappa\)。
\end{theorem}
\begin{proof}
将密度沿圆盘万有覆盖拉回，对圆盘使用 Ahlfors--Schwarz，再由 Poincaré 密度下降公式消去覆盖导数，给第一式。对 \(f\) 的拉回密度作同样比较给第二式，临界点照前一推论处理。
\end{proof}
这是以完备 Poincaré 度量为源的比较形式。完备性保障后续距离几何，但本证明不需要目标密度完备。若源只是任意完备度量，尚须额外的曲率下界与全局最大值工具；不能把本定理替换成“完备便足够”的泛化。
